\begin{table}[htbp]\centering
\begin{threeparttable}
\caption{The three naming systems. Each region of a copy is searched with BLAST+ against one external database per
system; every search requires an e-value of at most $10^{-50}$.}
\label{tab:arms}
\small
\begin{tabular}{lllp{0.30\linewidth}rr}
\toprule
System & Region & BLAST task & Database & Identity & Coverage\\
\midrule
VT  & SSU flank  & megablast & MaarjAM (gDAT, July 2020; 3{,}081 seq., 384 VT) & $\ge97\%$   & $\ge95\%$\\
SH  & full ITS   & blastn & UNITE 10.0, UNITE+INSD (16{,}800 Glomeromycota seq., 7{,}563 SH) & $\ge98.5\%$ & $\ge90\%$\\
LSU & LROR--FLR2 & blastn & Delavaux V18 (460 AMF records, 234 species labels) & $\ge97\%$ & $\ge90\%$\\
\bottomrule
\end{tabular}
\begin{tablenotes}\footnotesize
\item Coverage is of the shorter sequence, with reference Ns not counted; MaarjAM references span the NS31--AM1
fragment. Hits to any accession of the culture reference are dropped. The SH search was also run against the UNITE
general release (2{,}048 Glomeromycota representatives).
\end{tablenotes}
\end{threeparttable}
\end{table}
